\documentclass{article}
\usepackage{algorithm, algorithmic, amsfonts, amsmath}

\begin{document}
\begin{algorithm}[!h]
\renewcommand{\algorithmicrequire}{\textbf{Input:}}
\renewcommand{\algorithmicensure}{\textbf{Output:}}
	\caption{AlphaMax-N (class prior from noisy positives and unlabeled data)}
	\begin{algorithmic}[1]
	    \REQUIRE
	    Unlabeled sample $U=\{X_{i}^{u}\}$; \emph{noisy} positive sample $L=\{X_{i}^{l}\}$
	    (the labeled set itself may contain negatives)
	    \ENSURE
	    Class prior $\hat{\alpha}^{\ast}$ and noise level $\hat{\beta}^{\ast}$
	    \STATE \textbf{Setting:} unlike standard PU, the labeled positive set is contaminated,
	    so both $\alpha^{\ast}$ (prior) and $\beta^{\ast}$ (purity of $L$) must be recovered % weaker than standard PU: it does not even assume the labeled positive set is pure
	    \STATE \textbf{Step 0: univariate transform}
	    \STATE Fit a discriminative classifier separating $L$ from $U$ and map both samples through
	    its score % density estimation is infeasible in high dimension, so reduce to one dimension first
	    \STATE \quad \emph{Theorem 1: such transforms are $\alpha^{\ast}$-preserving, so the prior
	    is unchanged by this step} % this is why the reduction loses nothing of the target quantity
	    \STATE \textbf{Step 1: two applications of AlphaMax}
	    \STATE $\alpha^{+}\leftarrow \textsc{AlphaMax}(U,L)$
	    \hfill \emph{maximum proportion of $L$'s component in $U$}
	    \STATE $\beta^{+}\leftarrow \textsc{AlphaMax}(L,U)$
	    \hfill \emph{and symmetrically} % run once in each direction; this is the crux of the nonparametric route
	    \STATE \quad \textsc{AlphaMax}$(M,C)$ estimates the maximum proportion of component $C$
	    in mixture $M$ by maximising the combined likelihood over the mixing weight
	    \STATE \textbf{Step 2: recover the targets}
	    \STATE $(\hat{\alpha}^{\ast},\hat{\beta}^{\ast})\leftarrow$ apply Equation~(7) to
	    $(\alpha^{+},\beta^{+})$ % Lemma 1 and Theorem 1 give the closed-form relation between the two sets of quantities
		\RETURN $\hat{\alpha}^{\ast},\hat{\beta}^{\ast}$
	\end{algorithmic}
\end{algorithm}
% Jain, White, Radivojac, NeurIPS 2016 (arXiv 1606.08561), "Estimating the Class Prior
% and Posterior from Noisy Positives and Unlabeled Data". Compiled from Section 5,
% equation (7), Lemma 1 and Theorem 1 of the original.
% Difference from other prior estimators: TIcE, KM and BBE all assume the labeled positive set is pure. This method
% allows that set to be noisy, so it estimates two quantities at once, the class prior alpha* and the purity beta*.
% Two technical routes, given side by side in the original:
% (1) parametric MSGMM: model U and L as two Gaussian mixtures sharing components but differing in mixing proportions,
%     estimating (alpha, beta, u0, u1, Sigma0, Sigma1) jointly with an extended EM;
% (2) nonparametric AlphaMax-N, this pseudocode: call AlphaMax once in each direction, then invert in closed form.
% The one-dimensional transform is not a convenience. The original proves that such discriminative transforms preserve
% alpha*, so the curse of dimensionality is avoided without losing the quantity being estimated.
\end{document}
